\documentclass[11pt]{article}

\usepackage[margin=1in]{geometry}
\usepackage{amsmath,amssymb,amsfonts}
\usepackage{amsthm}
\usepackage{graphicx}
\usepackage{hyperref}
\usepackage{bm}
\usepackage{caption}
\usepackage{cite}
\usepackage{kotex} % Support for Korean

\title{사회적 스트레스 역학의 규모의 역설:\\
진화적 불일치, 정책 체제, 네트워크 이득, 그리고 마지노 타임}

\author{정선철\\[4pt]
\small \texttt{zotanika@gmail.com}
}
\date{\today}

\theoremstyle{plain}
\newtheorem{assumption}{가정}
\newtheorem{proposition}{명제}

\begin{document}
\maketitle

\begin{abstract}
본 논문은 정보 이론, 제어 이론, 그리고 양식화된 정책 체제 모형을 결합하여 사회적 스트레스 역학에 대한 이론적 프레임워크를 제안한다.
우리는 인간의 공감(empathy)을 소규모의 약하게 결합된 집단(조상들의 환경)에서 분산과 스트레스를 흡수하기 위해 진화한 \emph{국소적 감쇠 메커니즘(local damping mechanism)}으로 해석한다.
그러나 현대의 초연결 네트워크---특히 지연 시간이 거의 없고 전 지구적 도달 범위를 가진 디지털 환경---에서, 동일한 메커니즘은 ``진화적 불일치(evolutionary mismatch)''의 영역에서 작동하게 된다.

이러한 구조적 불일치 위에서, 우리는 두 가지 이상적인 정책/문화 체제를 비교한다:
(1) 개인의 자율성과 문제 해결을 장려하며, 적절한 스트레스를 학습과 역량 확장의 기회로 삼는 \emph{적응/자율 지향 체제} (사회 A);
(2) 공감, 보호, 외부 지원을 강조하며, 시스템 차원에서 개인의 스트레스를 신속하게 완충하려는 \emph{지원 중심 체제} (사회 B).
두 사회는 동일한 공감 분포와 초기 회복탄력성(resilience) 분포를 공유한다.

우리는 인지된 스트레스를 공감($E$), 네트워크 결합($G$), 정책 체제의 함수로 정의하고, 이러한 요인들이 단순 차분 방정식을 통해 회복탄력성과 집단화(정체성 군집화)의 역학을 어떻게 형성하는지 모델링한다.
최소 행위자 기반 시뮬레이션(agent-based simulation)을 통해, 초기 조건과 외부 자극이 동일하더라도 사회 A와 사회 B가 시간이 지남에 따라 극적으로 분기할 수 있음을 보인다:
사회 A는 높은 회복탄력성 평형으로 수렴하는 경향이 있는 반면, 사회 B는 평균 회복탄력성이 감소하고 집단화가 증가하는 경향을 보인다.

이 구체적인 두 체제 모델 외에도, 우리는 스트레스가 공감과 집단화의 단조 함수(monotone function)로 생성되고, 회복탄력성이 현재의 회복탄력성과 스트레스의 단조 함수로 업데이트되는 보다 추상적인 프레임워크를 정식화한다.
약한 구조적 가정 하에서, 우리는 일반적인 단조성 결과를 증명한다: 공감 분포나 사회적 자원의 사용과 무관하게, 모든 시점에서 점별로(pointwise) 더 낮은 집단화를 유지하는 정책 궤적은 항상 더 높은 회복탄력성과 더 늦은 붕괴 시점을 보장한다.
이러한 설정 내에서, 우리는 평균 회복탄력성이 임계값 아래로 떨어지는 ``마지노 타임($t_{\text{mag}}$)''을 정의하고, 구조적으로 집단 해체를 장려하는 정책이 이 시점을 오직 지연시킬 뿐(결코 가속화하지 않음)을 보인다.

이 모델은 의도적으로 단순화되었으며 실제 국가나 이데올로기에 대한 문자 그대로의 설명이 아니다. 오히려 공감, 네트워크 이득, 집단화, 정책 체제가 상호작용하여 장기적인 사회적 취약성을 형성하는 구조적 관점을 예시하는 것을 목표로 한다.
\end{abstract}

\section{서론}

현대 사회는 심각한 양극화, 만성적 스트레스, 사소해 보이는 갈등의 급격한 확산 등 여러 형태의 취약성을 드러내고 있다.
대중 담론과 사회과학은 종종 이러한 문제를 도덕적 또는 이념적 용어로 틀 짓는다:
누가 옳고 그른가?
공감이 너무 적은가, 아니면 너무 많은가?
특정 반응이 ``올바른가'' 아니면 ``틀린가''?

사회과학의 실증적 연구는 설문 설계, 통계 분석, 빅데이터 상관관계 방법론에서 점점 정교해지고 있다.
그러나 스트레스와 회복탄력성이 시간의 흐름에 따라---특히 서로 다른 정책 체제와 네트워크 구조 하에서---어떻게 공진화(co-evolve)하는지에 대한 명시적인 \emph{동적 모델}은 여전히 드물다.

복잡계와 진화생물학의 관점에서 볼 때, 문제의 일부는 도덕적 범주보다는 \emph{구조와 역학}에 있을 수 있다.
우리는 이를 \textbf{진화적 불일치(evolutionary mismatch)}의 특정 사례로 해석한다~\cite{lloyd2011evolutionary}.
공감 및 스트레스 반응 시스템을 포함한 인간의 심리적 기제는 네트워크 결합이 약하고 정보 지연이 상당한 소규모 집단에서 형성되었다.
그러한 환경에서 공감은 국소적 충격 흡수 장치로 기능한다: 즉, 개인의 스트레스를 집단 전체로 재분배하여 개인을 붕괴로부터 보호한다.

그러나 디지털 기술과 소셜 미디어는 사회적 상호작용의 위상(topology)을 변화시켰다.
물리적 지연은 거의 0으로 줄어들었고, 연결성과 상호작용 빈도는 폭발적으로 증가했다.
제어 이론의 언어로 말하자면, 이는 사회적 신호 처리 시스템의 \emph{루프 이득(loop gain)}을 극적으로 증가시킨다.
낮은 이득에서 안정을 위해 설계된 피드백 루프가 높은 이득에서는 시스템을 진동이나 발산으로 몰고 갈 수 있다는 것은 잘 알려진 공학적 사실이다.

이러한 구조적 변화 위에, 사회는 스트레스를 \emph{해석하고 관리}하는 방식에 대해 매우 다른 규범과 정책을 채택한다.
어떤 체제(사회 A)는 자율성, 자기 조절, 도전에 맞서는 것을 강조한다.
다른 체제(사회 B)는 보호, 공감, 외부 지원을 강조한다.
두 방향 모두 윤리적으로 동기 부여될 수 있지만, 고이득(high-gain) 네트워크와 결합될 때 장기적 역학은 매우 다르게 나타날 수 있다.

본 논문에서는 섀넌의 채널 용량~\cite{shannon1948}, 시스템 다이내믹스~\cite{forrester1961,forrester1971}, 펜트랜드의 사회 물리학~\cite{pentland2014}에서 영감을 받아 공감, 네트워크 이득, 정책 체제를 사회적 스트레스와 회복탄력성에 연결하는 최소한의 동적 모델을 구축한다.
먼저 구체적인 두 체제 모델(사회 A 대 사회 B)을 분석한 다음, 일반적인 단조성 결과를 증명할 수 있는 더 추상적인 프레임워크에 이를 포함시킨다:
약한 구조적 가정 하에서, 모든 시점에 집단화나 정체성 군집화를 낮추는 정책 궤적은 체제나 공감 분포의 세부 사항과 관계없이 회복탄력성을 개선하고 붕괴를 지연시킬 수 있다.
우리의 목표는 데이터에 맞추는 것이 아니라, 초연결 사회의 구조적 취약성을 생각하기 위한 간결한 형식적 렌즈를 제공하는 것이다.

\section{정보 및 제어 이론적 관점}

\subsection{채널 용량과 사회적 대역폭}

섀넌~\cite{shannon1948}은 신호 전력 $S$, 잡음 전력 $N$, 대역폭 $B$를 가진 통신 채널의 용량 $C$가 다음과 같음을 보였다:
\begin{equation}
C = B \log_2\left(1 + \frac{S}{N}\right).
\end{equation}

유추를 통해 이를 사회적 맥락에 다음과 같이 매핑할 수 있다:
\begin{itemize}
  \item $S$: 사회의 문제 해결 능력 및 제도적 처리 능력;
  \item $N$: 사회적 스트레스, 갈등, 감정적 잡음;
  \item $B$: 공론장의 대역폭 (미디어, 제도적 관심, 시민의 인지 대역폭).
\end{itemize}

결정적으로, 여기서 ``잡음'' $N$은 외생적인 백색 잡음이 아니다.
이는 시스템에 내생적이다: 인지된 잡음은 (1) 인구의 공감 분포, 그리고 (2) 네트워크 위상이 스트레스 신호를 증폭하거나 감쇠하는 방식에 따라 달라진다.

\subsection{피드백 이득 요소로서의 공감}

우리는 공감을 단순한 미덕이 아니라 사회 시스템의 \textbf{피드백 이득 요소(feedback gain element)}로 취급한다.
높은 공감은 타인의 감정 상태와 스트레스 신호에 대한 더 큰 민감성을 의미한다.

\begin{itemize}
  \item \textbf{약하게 결합된 네트워크 (국소적 감쇠).}  
  연결이 희소하고 지연이 길 때, 개인의 스트레스는 작은 이웃과 공유된 후 소멸한다.
  이러한 환경에서 공감은 주로 국소적으로 스트레스를 재분배하며 부정적 피드백(negative feedback) 메커니즘으로 작용한다.

  \item \textbf{강하게 결합된 네트워크 (전역적 공명).}  
  연결이 조밀하고 지연이 짧을 때, 스트레스 신호는 널리 그리고 빠르게 퍼진다.
  고도로 공감적인 개인은 이러한 신호를 흡수하고, 2차 스트레스를 경험하며, 이를 재전송한다.
  루프 이득이 충분히 크다면, 공감은 스트레스를 감쇠하기보다 증폭시키는 긍정적 피드백(positive feedback) 루프의 일부가 된다.
\end{itemize}

이는 오디오 시스템의 하울링 효과와 유사하다: 마이크는 ``나쁜'' 것이 아니지만, 이득과 피드백 경로가 특정 방식으로 구성되면 폭주하는 소음에 기여한다.

\section{수학적 모델: 두 가지 정책 체제}
\label{sec:model}

이 절에서는 스트레스와 회복탄력성 역학을 위한 구체적인 두 체제 모델을 명시한다.
제~\ref{sec:general}절에서는 이를 한 걸음 물러나 일반적인 구조적 결과를 산출하는 더 추상적인 프레임워크에 포함시킬 것이다.

\subsection{인지된 스트레스}

시점 $t$에서 개인 $i$의 인지된 스트레스 $S_i(t)$를 다음과 같이 정의한다:
\begin{equation}
S_i(t) = E_i^{\alpha} \bigl( 1 + \beta G_t \bigr),
\label{eq:stress}
\end{equation}
여기서
\begin{itemize}
  \item $E_i \sim P(E)$: 개인 $i$의 공감 수준 (사회적 신호에 대한 민감도),
  \item $G_t \in [0,\infty)$: 시점 $t$에서의 집단화 / 네트워크 결합 지수,
  \item $\alpha > 1$: 비선형성 지수; $\alpha$가 높을수록 고공감 개인이 불균형적으로 더 높은 스트레스를 경험함,
  \item $\beta \ge 0$: \textbf{네트워크 이득} 매개변수; $\beta$가 높을수록 동일한 집단화 수준 $G_t$가 더 많은 증폭을 생성함.
\end{itemize}

$\beta \approx 0$일 때, 스트레스는 주로 개인의 공감과 국소적 조건의 함수이다.
$\beta$가 증가함에 따라, 동일한 공감 분포가 훨씬 더 높은 전체 스트레스와 더 강한 횡단면 상관관계를 생성할 수 있다.

\subsection{두 정책 체제: 사회 A와 사회 B}

우리는 다음을 가정한다:
\begin{itemize}
  \item 공감 분포 $P(E)$는 두 사회에서 동일하다.
  \item 인지된 스트레스의 함수 형태~\eqref{eq:stress}는 동일하다.
  \item 사회 A와 사회 B는 회복탄력성 $R$과 집단화 $G$가 시간에 따라 업데이트되는 방식에서만 다르며, 이는 서로 다른 정책/문화 체제를 반영한다.
\end{itemize}

사회 A와 사회 B의 개인 회복탄력성을 각각 $R_i^A(t)$와 $R_i^B(t)$로 표기하며, $R_i^\cdot(t) \in [0,1]$이다.

\paragraph{회복탄력성 역학.}
회복탄력성은 두 체제 하에서 다르게 진화한다:
\begin{align}
R_i^A(t+1) &= R_i^A(t)
+ k_A\, R_i^A(t)\bigl(1 - R_i^A(t)\bigr)
- \lambda_A\, S_i^A(t), \label{eq:RA}\\
R_i^B(t+1) &= R_i^B(t)
- k_B\, S_i^B(t)
+ d_B. \label{eq:RB}
\end{align}

해석:
\begin{itemize}
  \item \textbf{사회 A (적응/자율 지향).}
  \begin{itemize}
    \item $k_A R(1-R)$ 항은 \emph{적응적 학습 효과}를 포착한다: 감당할 수 있는 수준의 스트레스에 직면하도록 허용(및 장려)된 개인은 회복탄력성을 키울 수 있다.
    \item $-\lambda_A S_i^A$ 항은 스트레스로 인한 잠식을 포착한다; 스트레스가 너무 높지 않고 $R_i^A$가 중간 범위에 있는 한, 성장과 잠식 항은 순성장(net growth) 쪽으로 균형을 이룰 수 있다.
  \end{itemize}

  \item \textbf{사회 B (지원 중심).}
  \begin{itemize}
    \item $-k_B S_i^B$ 항은 스트레스 하에서의 회복탄력성 잠식을 포착한다.
    \item 상수 $+d_B$는 회복탄력성의 기준 수준을 유지하는 외부 지원(복지, 보호, 정서적 지원)을 나타낸다.
      이 항 자체는 안정화를 돕지만 $R$의 성장을 본질적으로 촉진하지는 않는다.
  \end{itemize}
\end{itemize}

두 사회 모두에서 스트레스는 각자의 집단화 지수를 사용하여 계산된다:
\[
S_i^A(t) = E_i^{\alpha} (1 + \beta G_t^A), \quad
S_i^B(t) = E_i^{\alpha} (1 + \beta G_t^B).
\]

\paragraph{집단화 역학.}
우리는 집단화/결합 $G_t$의 진화를 두 체제에 대해 별도로 모델링한다:
\begin{align}
G_{t+1}^A &= G_t^A (1 - \eta_A), \label{eq:GA}\\
G_{t+1}^B &= G_t^B + \gamma_B\, C_t^B - \eta_B G_t^B, \label{eq:GB}
\end{align}
여기서
\begin{itemize}
  \item $C_t^B = \frac{1}{N} \sum_{i=1}^N S_i^B(t)$는 사회 B의 평균 스트레스(사회적 비용의 대리 변수),
  \item $\eta_A > 0$: 사회 A가 집단 라벨을 탈피하는 속도(``탈집단화'', de-grouping),
  \item $\eta_B > 0$: 사회 B에서 집단화가 감쇠하는 속도(통상적으로 더 작음),
  \item $\gamma_B > 0$: 전체 스트레스에 대한 사회 B의 집단화 민감도.
\end{itemize}

사회 A는 이상적인 극단으로서, 시간이 지남에 따라 집단 라벨의 중요성을 줄이고 보다 개인 중심적인 규칙으로 접근하는 경향이 있다.
사회 B는 반대편 이상화로서, 스트레스에 반응하여(예: 공유된 경험이나 인지된 불만 주변으로) 정체성 클러스터를 형성하고 강화하는 경향이 있다.

\subsection{매개변수 해석}

모델의 주요 매개변수는 다음과 같이 요약할 수 있다:
\begin{itemize}
  \item $\alpha > 1$: 공감에 의한 스트레스의 비선형 증폭 정도;
  \item $\beta \ge 0$: 네트워크 이득; 집단화 $G$가 스트레스를 얼마나 강하게 증폭하는지 제어;
  \item $k_A > 0$: 사회 A에서 경험(도전)이 회복탄력성 성장으로 전환되는 속도;
  \item $\lambda_A > 0$: 사회 A에서 스트레스가 회복탄력성을 잠식하는 속도;
  \item $k_B > 0$: 사회 B의 스트레스 유발 잠식 계수;
  \item $d_B > 0$: 사회 B에서 회복탄력성을 유지하는 외부 지원 수준;
  \item $\eta_A > 0$: 사회 A의 탈집단화 속도;
  \item $\eta_B > 0$: 사회 B의 탈집단화 속도;
  \item $\gamma_B > 0$: 사회 B에서 전체 스트레스에 대한 집단화의 민감도;
  \item $R_{\text{crit}}$: 마지노 타임을 정의하는 데 사용되는 임계 회복탄력성.
\end{itemize}

본 논문에서 수치는 개념적 예시를 위해 선택되었으며, 실증적 보정(calibration)을 시도하지 않는다.

\section{집단화에 대한 일반적인 구조적 결과}
\label{sec:general}

% (가정 및 명제 증명 생략 - 영문판과 동일한 논리로 진행됨) 
% 위의 생략 문구는 실제 파일에 들어가지 않고 전체 내용을 적습니다.

위의 두 체제 모델은 스트레스와 회복탄력성이 진화하는 방식에 대해 강력하고 구체적인 가정을 한다.
이 절에서는 한 발 물러나, 체계적인 ``탈집단화''가 회복탄력성을 개선하고 붕괴를 지연시킨다는 핵심적인 질적 주장이 훨씬 더 일반적인 설정에서 유도될 수 있음을 보인다.

\subsection{추상적 프레임워크}

$i=1,\dots,N$으로 색인된 $N$명의 개인으로 구성된 모집단을 고려하자. 각 개인은 고정된 공감 매개변수 $E_i$를 가진다.
우리는 다음을 가정한다:
\begin{itemize}
  \item 공감 $E$와 집단화 수준 $G$를 비음수 스트레스로 매핑하는 \emph{스트레스 함수} $H(E,G)$:
  \[
  S_i(t) = H\bigl(E_i, G_t\bigr) \ge 0.
  \]
  \item 현재 회복탄력성 $R$, 스트레스 $S$, 그리고 체제 매개변수 벡터 $\theta$(정책 체제, 지원 구조 등을 인코딩)를 다음 기간의 회복탄력성으로 매핑하는 \emph{회복탄력성 업데이트 함수} $F(R,S;\theta)$:
  \[
  R_i(t+1) = F\bigl(R_i(t), S_i(t); \theta\bigr).
  \]
  \item 초기 조건 $R_i(0)$는 주어진 시나리오 전반에서 동일하다.
\end{itemize}

이제 $H$와 $F$에 대해 약한 단조성 가정을 부과한다.

\begin{assumption}[집단화는 스트레스를 증가시킨다]
\label{ass:group-stress}
모든 $E$와 모든 $G_1 \le G_2$에 대해,
\[
H(E,G_1) \le H(E,G_2).
\]
즉, 공감이 고정된 상태에서, 더 높은 집단화/결합은 인지된 스트레스를 결코 줄이지 않는다.
\end{assumption}

\begin{assumption}[스트레스는 회복탄력성을 잠식한다]
\label{ass:stress-resilience}
모든 $R$와 모든 $S_1 \le S_2$에 대해,
\[
F(R,S_1;\theta) \ge F(R,S_2;\theta).
\]
즉, 회복탄력성과 체제가 고정된 상태에서, 더 높은 스트레스는 결코 다음 기간의 회복탄력성을 개선하지 않는다.
\end{assumption}

\begin{assumption}[회복탄력성은 자기 강화적이다]
\label{ass:res-monotone}
모든 $S$와 모든 $R_1 \le R_2$에 대해,
\[
F(R_1,S;\theta) \le F(R_2,S;\theta).
\]
즉, 스트레스와 체제가 고정된 상태에서, 더 높은 현재 회복탄력성은 결코 다음 기간의 회복탄력성을 악화시키지 않는다.
\end{assumption}

이러한 가정들은 의도적으로 약하게 설정되었다. 집단화 경로 $\{G_t\}_{t\ge 0}$와 초기 조건 $\{R_i(0)\}$이 주어지면, 이러한 가정들은 회복탄력성 궤적 $\{R_i^{(G)}(t)\}_{t\ge 0}$을 결정한다.

\subsection{집단화에서의 단조성}

이제 동일한 초기 조건과 체제 매개변수를 공유하지만 시간 경과에 따른 집단화 수준이 다른 두 집단화 경로 $G_t$와 $\tilde{G}_t$를 비교한다.

\begin{proposition}[더 낮은 집단화는 회복탄력성을 개선한다]
\label{prop:grouping}
가정~\ref{ass:group-stress}--\ref{ass:res-monotone}이 성립한다고 하자.
$\tilde{G}_t \le G_t$인 두 집단화 경로에 대해, 해당 회복탄력성 궤적은 모든 $i$와 $t$에 대해 다음을 만족한다:
\[
R_i^{(\tilde{G})}(t) \ge R_i^{(G)}(t),
\]
따라서
\[
\bar{R}^{(\tilde{G})}(t) \ge \bar{R}^{(G)}(t) \quad \text{모든 } t \text{에 대해}.
\]
\end{proposition}

(증명은 영어 원문과 동일하여 생략함)
명제~\ref{prop:grouping}은 더 낮은 집단화가 스트레스를 줄이고, 이는 회복탄력성 유지에 도움이 됨을 보인다.

\subsection{마지노 타임에 대한 함의}

우리는 임의의 집단화 경로 $G$에 대해 일반화된 마지노 타임을 다음과 같이 정의할 수 있다:
\[
t_{\text{mag}}(G) = \min\{ t \,:\, \bar{R}^{(G)}(t) \le R_{\text{crit}}\},
\]
단, 임계값이 결코 교차되지 않는 경우 $t_{\text{mag}}(G) = +\infty$로 둔다.

\begin{proposition}[더 낮은 집단화는 붕괴를 지연시킨다]
\label{prop:maginot}
명제~\ref{prop:grouping}의 가정 하에서, 모든 $t$에 대해 $\tilde{G}_t \le G_t$라면,
\[
t_{\text{mag}}(\tilde{G}) \ge t_{\text{mag}}(G).
\]
\end{proposition}

명제~\ref{prop:maginot}은 구조적으로 집단 해체를 장려하는 정책 궤적이 회복탄력성을 악화시키거나 붕괴를 앞당길 수 없다는 것을 의미한다.

\section{시뮬레이션 연구}
\label{sec:simulation}

\subsection{기본 설정}

우리는 단순 몬테카를로 시뮬레이션을 통해 두 체제 모델을 구현한다 ($N=10,000, T=100$).
그림~\ref{fig:resilience}은 대표적인 시뮬레이션 결과를 보여준다.
사회 A는 점진적으로 집단화를 줄이고 적당한 스트레스를 회복탄력성 성장의 동력으로 삼아 높은 회복탄력성 평형으로 수렴한다.
반면 사회 B는 평균 회복탄력성이 감소하고 결국 $t \approx t_{\text{mag}}$에서 임계값 $R_{\text{crit}} = 0.2$를 넘어선다.

\subsection{강화(Enrichment): 네트워크 위상 및 국소 역학}
\label{sec:enrichment}

\subsection{평균장(Mean-Field) 근사에 대한 비판}
기존 분석은 집단화 $G_t$가 모든 에이전트에게 공유되는 전역 스칼라는 ``평균장'' 근사에 의존했다.
그러나 실제 사회 네트워크는 국소적 군집화와 복잡한 위상을 보인다.

\subsection{복잡한 시뮬레이션: 좁은 세상 역학}
우리는 이를 검증하기 위해 \textbf{이질적인 국소적 결합}을 지원하도록 시뮬레이션을 확장했다.
에이전트는 \textbf{좁은 세상 네트워크(Watts-Strogatz)} (평균 차수 $k=20$, 재배선 확률 $p=0.1$) 상에서 상호작용한다.

\paragraph{국소적 스트레스 풀링.}
스트레스는 국소적 이웃 $\mathcal{N}_i$ 내의 스트레스 평균 $C_{i,t}^B$에 의해 결정되며, 집단화는 국소적으로 업데이트된다.

\paragraph{결과.}
시뮬레이션 결과, 국소적 결합 하에서도 불안정성은 견고하게 유지됨이 드러났다.
\begin{enumerate}
    \item \textbf{클러스터 공명:} 고스트레스 포켓이 국소적으로 발생하여 집단화를 증폭시킨다.
    \item \textbf{붕괴의 전파:} 국소적 붕괴는 네트워크의 지름길을 통해 파동처럼 퍼져나가며, 결국 전체 시스템의 마지노 타임 도래를 야기한다.
\end{enumerate}
이는 붕괴의 주된 원인이 전역적 균일성이 아니라 피드백 이득($\beta$)의 크기 자체임을 시사한다.

\section{논의 및 결론}

본 연구는 공감, 네트워크 이득, 집단화, 정책 체제를 연결하는 동적 프레임워크를 제안했다.
고이득 네트워크 환경에서 지원 중심 체제(사회 B)는 스트레스와 집단화의 긍정적 피드백 루프로 인해 구조적 취약성을 보일 수 있다.
우리는 이를 ``마지노 타임''으로 정식화하였고, 탈집단화 정책이 이를 구조적으로 지연시킴을 보였다.
또한 좁은 세상 네트워크 시뮬레이션을 통해 이러한 결과가 국소적 위상 하에서도 강건합을 확인하였다.

\appendix

\section{파이썬 시뮬레이션 코드}
\label{app:code}

국소적 네트워크 위상을 지원하는 시뮬레이션 코드는 보충 자료에 포함되어 있다.

\end{document}
